\section{Introduction}
\label{sec:intro}

Deep learning (DL) programs are written, reviewed, and debugged as Python code, but they are increasingly
executed in another form. Since PyTorch~2, a call to \tcompile routes a model through bytecode-level graph
capture, functionalization and automatic-differentiation tracing, operator decomposition, loop-level lowering,
C++ or Triton code generation, guard-based specialization, and several caches~\cite{DBLP:conf/asplos/AnselYHGJVBBBBC24}.
Comparable pipelines exist in TVM~\cite{DBLP:conf/osdi/ChenMJZYSCWHCGK18}, XLA-backed TensorFlow~\cite{DBLP:conf/osdi/AbadiBCCDDDGIIK16}, MLIR-based
stacks~\cite{DBLP:conf/cgo/LattnerAB0DPRSV21}, and Triton~\cite{DBLP:conf/pldi/TilletKC19}; we call such a pipeline an \emph{AI compiler stack}. These systems are expected to produce an optimized artifact
that is semantically equivalent to eager execution, so that the Python source remains the single description of
what the program does.

This expectation can be violated without any visible sign in the source. An empirical study of \emph{toxic
compilation}\footnote{Reference withheld for double-anonymous review.} collected more than a thousand incidents
in which the produced artifact diverged from the reviewed source. We call their common outcome a
\emph{source--artifact consistency failure}: for the same program, input, and context, eager and compiled
execution differ in an observable way. Four of the defects reported in this paper illustrate the range. An
\code{OrderedDict.move\_to\_end} call inside a compiled function is dropped, so that an LRU cache maintained
next to the compiled step evicts the entry it has just used. An \code{IndexError} raised by a tensor operator
bypasses the \code{except} clause written for it. A generator that leaves the compiled region is returned as an
exhausted tuple iterator whose return value is \code{None}. The expression \code{x * 1} at the end of a graph
returns the input tensor itself, so that an in-place update of the result modifies the input.

Compiler testing has so far varied the \emph{program}. NNSmith~\cite{DBLP:conf/asplos/LiuLRTLPZ23}, NeuRI~\cite{DBLP:conf/sigsoft/0004PW023},
and HirGen~\cite{DBLP:conf/issta/MaS00C23} synthesize computation graphs; Tzer~\cite{DBLP:journals/pacmpl/LiuWYDZ22} and
MLIRSmith~\cite{DBLP:conf/kbse/WangCXLWSZ23} mutate intermediate representations; TorchProbe~\cite{DBLP:conf/aplas/SuGPS23} applies
semantics-preserving transformations to Python programs; FreeFuzz~\cite{DBLP:conf/icse/WeiDYZ22},
DeepREL~\cite{DBLP:conf/sigsoft/DengYWZ22}, and LLM-based fuzzers~\cite{DBLP:conf/issta/DengXPY023,DBLP:conf/icse/DengXYZY024} generate API
invocations. In all of these techniques a program is compiled along one path, the one the harness takes by
default. For \tcompile this leaves much of the compiler untested. Where TorchDynamo breaks the graph depends on
the call site; whether AOTAutograd participates depends on whether a gradient is needed; whether TorchInductor
emits a scalar or a vectorized loop depends on the tensor length; whether an operator takes a dedicated lowering
or a generic decomposition depends on its dtype; a \code{TorchDispatchMode}, autocast, or CUDA graphs change
which components run; export and ahead-of-time packaging replace the runtime. A defect on a path the harness
does not take remains undetected regardless of how many programs are generated.

The compilation itself can, however, be varied: its path can be enumerated and controlled one factor at a time. The statements after
which a graph break can be inserted, the backend layers that can be selected, the execution contexts that can
be entered, and the artifact routes that can be taken are switches exposed by the pipeline; their number is
small, their meaning is documented, and none of them is permitted to change the result of the program. The
program-side factors that change compiler behavior can be enumerated as well: a static analysis of the source
recovers the shapes, dtypes, scalars, and flags a compilation site reads, and operator metadata lists the dtypes
and sample shapes each operator supports. Treating both groups as test dimensions, holding one fixed and
changing the other by one factor at a time, yields tests in which every difference has a single candidate cause.
A path factor is a semantics-preserving change to how the program is compiled, so the eager result is the
expected outcome of a metamorphic relation~\cite{DBLP:journals/csur/ChenKLPTTZ18}; the relation is the same for
every factor, and the design effort goes into the choice of observables that are compared.

We implement this approach in \tool. Section~\ref{sec:approach} defines the factor model, the one-factor
execution matrix, an oracle over seven observables (values against a float64 or multi-precision reference,
dtype and shape and stride, object identity and aliasing, exception type and its routing to a handler, side
effects on inputs and external state, gradients, and process survival), and layered localization, which names
the stage at which a difference first appears. Section~\ref{sec:analyzers} describes how the path factors are
realized on \tcompile and how program-side factors are drawn from the source and from operator metadata.

\paragraph{Results.} Reports filed against PyTorch~2.14 and nightly builds during the study number 73: 55 new
issues, 15 comments that add members or root causes to existing issues, and 3 reports to the vendor of the C++
compiler that TorchInductor uses on Windows. At the time of writing, 14 are fixed or have a fix under review, 10
are confirmed by maintainers or reproduced by other contributors, 48 await triage, and 1 was judged expected
behavior. A corpus of 185 programs that program-only differential testing under three configurations had already
covered yielded two new Dynamo defects when one path factor, the position of graph breaks, was added. One
metadata-derived factor exposed 134 operators that compute under \tcompile on dtypes their eager kernels
reject. Of the 58 findings attributable to the compile pipeline, 39 are detectable only through an observable
other than the output value. Guards, caches, compile composition, autocast, higher-order operators, custom
operators, and the numerical behavior of the export path were consistent in every configuration we tried.

\paragraph{Contributions.}
\begin{itemize}
  \item Varying the compilation as a dimension of consistency testing, with a set of path factors that a
        harness can enumerate and control (graph-break position, backend layer, execution context, artifact
        route), and evidence that program-only testing does not reach the defects on rarely taken paths
        (Sections~\ref{sec:approach}, \ref{sec:analyzers}, \ref{sec:rq1}).
  \item A one-factor execution matrix over program-side and path-side factors and an oracle over seven
        observables, which make every difference attributable and make side-effect, identity,
        exception-routing, and process-level failures detectable (Sections~\ref{sec:approach}, \ref{sec:rq3}).
  \item An evaluation on injected faults, historical issues, and current PyTorch builds, with 73 reports
        classified by maintainer response and by root-cause location, and negative results that state where
        consistency holds (Section~\ref{sec:eval}).
\end{itemize}
